% sections/erratum.tex -- erratum to version 1 (arXiv:2508.14429v1).
% Intended to be \input right after sections/theory.tex: it is the last
% subsection of the theory section (change \subsection to \section if it is
% placed elsewhere). Every instance below is produced and checked by
% research/mmhm/reference/check_props.py (recorded in reference/README.md and
% reference/counterexamples.json).

\subsection{Erratum to version 1}\label{sec:erratum}

This paper replaces version~1~\cite{ouyang2025mmhm_v1}. The update rule of
version~1 does not maintain exact Betti numbers. We state what version~1
specified, give counterexamples to its two exactness claims, and record an
inconsistency in its experiment P2. Section and figure numbers below refer to
version~1.

\paragraph{What version 1 specified.}
An acyclic matching on $K$ defines critical cells, critical chain groups
$C_k^C$ and a critical boundary $\bd_k^C$ ``induced by the matching via
$V$-paths (Forman's gradient flow)'', with matrices $\Bm{k}^C$; Betti numbers
are read off from $r_k=\rank\Bm{k}^C$ and the numbers of critical cells
(\S2.2, \S2.5). The matching is computed at initialisation and at
\emph{recompression events}: every $m$ updates, when the fraction $\rho_t$ of
affected critical cells reaches a threshold $\tau$, or when a gate precondition
fails (typical choices $m=32$, $\tau=0.30$). No rule is given for updating the
matching between recompressions. For an edit, $R_k\subseteq\Kc{k}$ is the set of
$k$-simplices ``whose incidences in the boundary operators are affected by the
modification'', computed by comparing the stars of the modified region before
and after the update, and the affected critical columns are
\[
  R_k^C=\{\,j : \sigma_j^k\in C_k^C\ \text{and}\ \sigma_j^k\in R_k\,\}
\]
(\S2.3). The columns $\{\Bm{k}^C[:,j] : j\in R_k^C\}$ are replaced by their new
incidence patterns, ``followed by reapplying column-reduction operations only to
these affected columns and to any columns whose pivots conflict with them'';
``when a column is replaced, its pivot assignment is cleared, the column is
re-reduced against existing pivots, and any resulting pivot conflicts are
propagated in a localized manner''. If the row set of $\Bm{k}^C$ changes, the
pivot map of that dimension is re-initialised. Version~1 states that this
``ensures that the rank of $\Bm{k}^C$ is updated exactly'' (\S2.3) and
``guarantees exact Betti numbers'' (\S2.5). Both statements are false, for two
independent reasons, E1 and E2.

\paragraph{E1: a reduced column records the columns it absorbed.}
In a column reduction $R=BV$, a reduced column $R_b=\bd b+\bd a_1+\dots$ is
correct only while the absorbed columns $a_i$ are present with their old
boundaries. The rule above re-reduces a replaced column $a$, but leaves every
$R_b$ that absorbed $a$ in place. Such an $R_b$ need not own a pivot (it may be
zero), so it is never in conflict with anything, and nothing in the rule
revisits it. A deletion is a replacement by the zero column and fails in the
same way. The DRB repairs exactly this: $\occ(a)$ lists every slot whose $V$
contains $a$, and $\Del$ rewrites all of them (Theorem~\ref{thm:exact}(c)).

\begin{example}[E1, matrix form]\label{ex:e1-matrix}
Rows $e_1<e_2$, $B=[\,a\mid b\,]$ with $a=b=e_1$. Reduction: $R_a=e_1$ owns pivot
$e_1$; $R_b=b+a=0$; rank $1$. Replace $a$ by $e_2$: its pivot is cleared, it is
re-reduced to $R_a=e_2$, which owns the free pivot $e_2$, so there is no
conflict; $R_b=0$ is not touched. The reported rank is $1$; the rank of
$[\,e_2\mid e_1\,]$ is $2$.
\end{example}

\begin{example}[E1, shortest simplicial instance]\label{ex:e1-simplicial}
The empty matching is admissible in version~1 (every simplex in at most one pair,
no closed $V$-path); then $\Bm{k}^C=\Bm{k}$. Apply, one edit per update, with keys
in insertion order,
\[
  +\{0\},\ +\{1\},\ +\{2\},\ +\{0,1\},\ +\{0,2\},\ +\{1,2\},\ -\{0,1\}.
\]
After $+\{1,2\}$ the reduction gives $R_{01}=\{0,1\}$, $R_{02}=\{0,2\}$ and
$R_{12}=\{1,2\}+\{0,2\}+\{0,1\}=0$, so $\beta=(1,1)$. Deleting $\{0,1\}$ clears
its pivot and drops its column; $R_{12}=0$ stays. Reported: $r_1=1$,
$\beta=(2,1)$. True: the path $1$--$2$--$0$, $r_1=2$, $\beta=(1,0)$. The DRB has
$V_{12}=\{\{1,2\},\{0,2\},\{0,1\}\}$, so $\occ(\{0,1\})$ contains the slots of
$\{0,1\}$ and $\{1,2\}$; Case A applies with $z$ the slot of $\{1,2\}$, sets
$V_{01}\gets\{\{1,2\},\{0,2\}\}$ (still $\bd V_{01}=\{0,1\}=R_{01}$), removes
slot $z$, and reports $r_1=2$.
This is a shortest instance. An exhaustive enumeration of the valid edit
sequences on at most four vertices, up to relabelling (1, 2, 5, 15, 61 and 263
sequences of lengths 1 to 6, and 1250 of length 7), finds no failure of length
at most 6 and 18 of length 7; a sequence of length at most 6 on five or more
vertices contains at most one edit of positive dimension and cannot fail. The failing
edit deletes an edge, so the row set of $\Bm{1}$ is unchanged and the
re-initialisation clause does not apply. Adding $N\ge5$ isolated vertices, which
are critical under every matching and untouched by the edits, makes
$\rho_1\le 3/(6+N)<0.3$; with $m=32$ and gates disabled (as in version~1's P2),
no recompression is triggered, and the reported $\beta_0$ is one too large and
$\beta_1=1$ instead of~$0$.
\end{example}

\paragraph{E2: the affected set misses distant critical columns.}
The critical boundary is a sum over gradient paths, $\dM c=\flow(\bd c)$
(Section~\ref{sec:morse}), and it depends on every pair along every gradient
path leaving the facets of $c$. These paths can be arbitrarily long, so an edit
to a matched cell can change $\dM c$ for critical cells $c$ far from the edit,
while $R_k^C$ contains only critical cells whose own incidences change. Since
version~1 does not say how the matching is maintained between recompressions, we
use the least intrusive convention: keep every pair that does not involve a
deleted cell; the partner of a deleted matched cell, and every inserted cell,
becomes critical. The restriction of an acyclic matching is acyclic, so this is
a valid matching. (Under this convention insertions are harmless: a new cell is
maximal and critical and lies on no gradient path. Deleting a matched cell is
not.)

\begin{example}[E2 on cycle graphs]\label{ex:e2-cycle}
Let $K=C_L$, $L\ge4$, with the collapse matching of
Definition~\ref{def:collapse}. The collapse removes one critical edge $c$, then
collapses the remaining path completely, leaving one critical vertex $r$. The
gradient paths from the two endpoints of $c$ are the two arcs of $C_L\setminus c$
from those endpoints to $r$; together they contain every matched edge exactly
once, and $\dM c=r+r=0$. Delete a matched edge $y=V(x)$: now $x$ is critical,
the arc through $y$ ends at $x$, and $\dM c=r+x\ne0$. When $y$ shares no vertex
with $c$, the column $c$ is not in $R_1^C$ and stays $0$; the critical complex
then reports $\beta=(2,1)$ instead of $(1,0)$. Such $y$ exist at graph distance
up to about $L/2$ from $c$. The smallest case, $L=4$ with edges
$\{0,1\},\{1,2\},\{2,3\},\{0,3\}$: the collapse gives $c=\{0,1\}$, $r=\{3\}$ and
pairs $(\{1\},\{1,2\})$, $(\{0\},\{0,3\})$, $(\{2\},\{2,3\})$; deleting $\{2,3\}$
changes $\dM\{0,1\}$ from $0$ to $\{2\}+\{3\}$.
\end{example}

\begin{example}[E2 with a critical triangle]\label{ex:e2-torus}
Let $K$ be the $3\times3$ grid torus on the vertices $1,\dots,9$, in which
vertex $(i,j)$, $i,j\in\mathbb{Z}/3$, is labelled $1+3i+j$ and the triangles are
\[
  \{(i,j),(i{+}1,j),(i{+}1,j{+}1)\},\qquad \{(i,j),(i,j{+}1),(i{+}1,j{+}1)\},
\]
together with the cone with apex $0$ over the meridian $1$--$4$--$7$;
$\beta(K)=(1,1,1)$. The collapse matching (FIFO queue) has critical cells
$\{0,1,4\}$, $\{1,2,5\}$, $\{3,4\}$, $\{3,8\}$, $\{9\}$, and $\dM\{1,2,5\}=0$.
Deleting the matched triangle $\{6,7,9\}$, paired with the edge $\{6,9\}$,
changes $\dM\{1,2,5\}$ to $\{6,9\}$, although $\{1,2,5\}$ and $\{6,7,9\}$ share no
vertex. The stale critical complex reports $\beta=(1,2,1)$; the truth is
$(1,1,0)$.
\end{example}

\paragraph{P2: an inconsistent workload description.}
Version~1, \S3.3, defines the P2 sequence as alternating a refinement and ``its
exact inverse'', $\textsc{Refine}(R_t)$ followed by $\textsc{Coarsen}(R_t)$. Under
that definition $|K_t|$ takes at most two values. The same section states that
each step injects $O(1)$ new simplices, ``thus $|C_t|=\Theta(t)$'', and
Figures~4 and~5 of version~1 show sizes growing linearly over 500 updates. The
two descriptions cannot hold for the same sequence, and version~1 does not
determine which sequence produced its P2 timings. The workload W2 of this paper
lets the size perform a random walk around its initial value and records it.

\paragraph{Why the experiments of version 1 did not reveal E1 and E2.}
P1 (an octahedron with a two-triangle window) and P3 (a connected genus-$0$
shell with ports) are connected surfaces with known connectivity, whose Betti
numbers follow from $\chi$, $\beta_0=1$ and the boundary-edge criterion without
any matrix (Remark~\ref{rem:surfaces}); version~1's gates (\S2.4) are these
formulas. P2 has the constant Betti numbers $(1,0,0)$. None of the three
experiments can therefore distinguish a correct update from an incorrect one.
In this paper every exact method is compared with independent recomputations
(Section~\ref{sec:exp-correctness}), also on three-dimensional and
non-manifold complexes whose Betti numbers change.

\paragraph{Cost claims.}
Version~1, \S2.6, states a local reduction cost of
$\Theta\bigl(\sum_k (s_k^C)^2\bigr)$ ``under sparse operations with limited
fill-in''. That analysis concerns the rule shown above to be inexact, and we
make no statement about it. For the exact method of this paper,
Proposition~\ref{prop:nonlocal} shows that no bound in terms of the size of the
update holds, and we report costs as measurements.
